本次研究构建了支持\textbf{重叠/不重叠纤维生成}与\textbf{薄膜桥状binder建模}的GDL三维微结构生成方法，完成四种模式的微结构生成且孔隙率均逼近目标值，验证了方法的可行性，同时明确了当前模型存在的局限性与后续优化方向。

\subsection{研究的改进与创新特征}
本研究在微结构建模方法上实现多项改进与创新：提出厚度分层与高斯抖动采样法，突破碳棒纤维随机分布假设，精准复现 GDL 层状堆叠的结构特征；融合空间分箱、AABB 预筛与参数自适应收缩策略，有效解决不重叠模式下的 jamming 问题，提升模型生成稳定性；建立纤维接触几何驱动的 binder 靶向建模法，摒弃传统均匀填充思路，实现胶膜与胶桥的差异化形态构建；设计二值体素场软边界平滑算法，缓解体素化带来的锯齿伪影，显著优化微结构可视化效果；开发多模式切换建模框架并制定标准化输出流程，突破单一建模策略局限，有效提升方法的通用性与复用性。

\subsection{研究存在的局限性}
本研究仍存在若干不足之处。binder 形态表征的真实性有待提升，不重叠模式下生成的 binder 多以异面胶桥为主，而在重叠模式下局部区域易出现凸起特征，与实际结构分布存在偏差；不重叠模式（opt1/opt3）仍存在多次采样失败与 jamming 问题，算法整体生成效率仍有提升空间。现有平滑算法对局部细节的处理能力有限，在细小 binder 胶桥、纤维端部等区域仍存在较为明显的锯齿伪影；不重叠模式下生成的碳纤维中短细纤维占比偏高，且整体分布偏规整，与实际 GDL 碳纤维随机堆叠的特征存在差异。同时，模型将碳纤维简化为等半径理想长直圆柱体，未考虑纤维弯曲、粗细不均等真实几何特征，一定程度上影响了接触几何表征的准确性。此外，当前建模框架尚未引入 binder 亲疏水性、热导率等关键物理参数，暂难以实现 GDL 结构–物理–性能的多场耦合分析。

\subsection{后续优化方向}
后续将围绕微结构建模与仿真开展系统性优化：通过优化 binder 形态分类规则与厚度函数参数、提升纤维共面接触对占比并强化凹弧控制以修正局部凸起问题；引入 KD 树空间索引与多参数分级收缩策略，降低非重叠模式下碰撞检测计算开销并减少 jamming 现象；采用几何轮廓引导的自适应平滑算法，对局部细节区域精细化调控平滑核以提升界面平滑度；构建纤维半径分布与弯曲几何模型并引入随机粗糙度扰动，使碳纤维表征更贴近真实结构特征；建立 binder 物理性质赋值机制，明确几何形态与物理参数的关联关系，夯实多物理场建模基础；同时结合实验表征数据校准模型输入参数，通过定量指标验证微结构真实性，最终形成实验与仿真相互支撑、迭代完善的闭环优化体系。

\subsection{结论}
本次研究通过厚度分层采样、碰撞检测加速、binder靶向生成等策略，实现了GDL纤维相、binder相、孔隙相的三相微结构耦合表征，在层状结构复现、binder形态精准建模等方面的改进具有实际应用意义。但本方法在形态真实性、算法生成效率、物理性质表征等方面仍存在局限性，后续将通过模型精细化、算法优化与实验数据校准进一步提升方法性能，为GDL结构-性能耦合分析提供更可靠的几何模型支撑。